\section{Lacunas Econômicas: A Viabilidade do Modelo de Negócio}

A viabilidade econômica dos gêmeos digitais odontológicos é talvez a dimensão menos explorada na literatura de cunho estritamente técnico, mas potencialmente a mais decisiva para a sua adoção em larga escala. A sustentabilidade financeira da tecnologia não se restringe ao preço do software, abarcando toda a cadeia de valor.

\subsection{A Ausência de Estudos de Custo-Efetividade}

Estudos formais de custo-efetividade (cost-effectiveness analysis) e de custo-utilidade (cost-utility analysis) envolvendo gêmeos digitais odontológicos são virtualmente inexistentes na base Scopus ou PubMed. As perguntas pragmáticas de gestores clínicos permanecem sem resposta documentada: o alto investimento financeiro inicial (Capex) e operacional (Opex) em infraestrutura de TI, licenças corporativas e treinamento de equipe se traduz em uma redução percentual de falhas e retratamentos suficiente para justificar o desembolso (\textit{Return on Investment - ROI})?

Para o Sistema Único de Saúde (SUS), onde o orçamento \textit{per capita} para a saúde bucal é severamente limitado e disputado (\autoref{ch:impacto}), a equação é ainda mais draconiana. Um planejamento que onera em R\$ 200 adicionais um tratamento reabilitador só é justificável, sob a ótica da avaliação de tecnologias em saúde (ATS), se induzir a uma redução de desfechos negativos ou iatrogenias clinicamente relevante. A ausência de dados longitudinais de larga escala inviabiliza que a Comissão Nacional de Incorporação de Tecnologias no SUS (CONITEC) produza pareceres favoráveis de incorporação.

\subsection{Modelos de Financiamento Fragmentados}

Além da mensuração de eficiência, persiste o impasse do modelo de financiamento. Na odontologia privada, o custo é invariavelmente repassado ao paciente final (Out-of-Pocket), confinando o gêmeo digital a um nicho de altíssima renda (\textit{Premium Care}). No sistema de saúde suplementar, as operadoras de planos odontológicos relutam veementemente em codificar procedimentos baseados em IA sem evidências de redução de sinistralidade de longo prazo. O resultado é um ciclo vicioso de estagnação: sem reembolso, não há escala; sem escala, não há coleta de dados do mundo real (\textit{Real-World Evidence}); e, sem evidências, não há aprovação de reembolso.
